\documentclass[11pt]{article}

\usepackage[margin=1in]{geometry}
\usepackage{amsmath, amssymb}
\usepackage{xcolor}
\usepackage{listings}
\usepackage{hyperref}
\usepackage{array}
\usepackage{booktabs}

\hypersetup{
    colorlinks=true,
    linkcolor=blue!60!black,
    urlcolor=blue!60!black,
    citecolor=blue!60!black
}

\lstdefinestyle{bugpython}{
  language=Python,
  basicstyle=\ttfamily\small,
  keywordstyle=\color{blue!60!black},
  stringstyle=\color{green!40!black},
  commentstyle=\color{gray!70!black},
  showstringspaces=false,
  breaklines=true,
  frame=single,
  rulecolor=\color{gray!30},
  columns=fullflexible,
  keepspaces=true
}

\title{SymPy Bug Report}
\author{}
\date{}

\begin{document}

\maketitle

\section{Bug 10: Principal Cube-Root Inequality Includes Negative Reals}

\subsection{Minimal Reproducer}

\medskip

\begin{lstlisting}[style=bugpython]
from sympy import S, N, symbols
from sympy.solvers.inequalities import solve_univariate_inequality

x = symbols("x")
sol = solve_univariate_inequality(x**(S(1)/3) <= 1, x, relational=False)

print("solution:", sol)
print("contains_minus_one:", sol.contains(-1))
print("lhs_at_minus_one:", N((-1)**(S(1)/3), 30))
\end{lstlisting}

\subsection{Observed Behavior}

\medskip

\begin{lstlisting}[style=bugpython]
solution: Interval(-oo, 1)
contains_minus_one: True
lhs_at_minus_one: 0.5 + 0.866025403784438646763723170753*I
\end{lstlisting}

SymPy reports that the solution set is \((-\infty, 1]\). In particular, the
reported set contains \(x=-1\), even though SymPy's principal value of
\((-1)^{1/3}\) is non-real. Such a point cannot satisfy a real-valued inequality
of the form \(x^{1/3} \leq 1\).

\subsection{Expected Behavior}

The mathematically correct solution over the real line is
\[
    [0, 1].
\]
More generally, for the parametrized family checked in the artifact,
\[
    \operatorname{solve\_univariate\_inequality}
    \left(x^{1/3} \leq m,\ x,\ \texttt{relational=False}\right)
    = [0, m^3]
\]
for the positive integer values \(m=1,\ldots,6\).

\subsection{Mathematical Explanation}

SymPy's expression \(x^{1/3}\), represented as \verb|x**(S(1)/3)|, uses the
principal branch of the power function. For a negative real input \(x=-a\) with
\(a>0\),
\[
    (-a)^{1/3}
      = \exp\!\left(\frac{\log(a)+i\pi}{3}\right)
      = a^{1/3}\left(\frac{1}{2} + \frac{\sqrt{3}}{2}i\right),
\]
which is not real. Therefore the real-valued domain of the left-hand side is not
all of \(\mathbb{R}\); for this inequality it is \(x \geq 0\).

On that domain, the principal cube root agrees with the ordinary nonnegative
real cube root and is monotone increasing. Thus
\[
    x^{1/3} \leq 1
    \quad\Longleftrightarrow\quad
    0 \leq x \leq 1.
\]
The returned interval \((-\infty,1]\) is not an equivalent representation or a
harmless branch convention: it contains points at which the left-hand side is
complex, so the ordered real comparison is not satisfied.

\subsection{Source-Level Diagnosis}

The diagnosis identifies the primary source-level issue in
\nolinkurl{sympy/calculus/util.py}, in the function \texttt{continuous\_domain}.
The public call enters \texttt{solve\_univariate\_inequality} in
\nolinkurl{sympy/solvers/inequalities.py}. For the representative input, the
solver forms \(e=x^{1/3}-1\), obtains the equality root \(x=1\), and asks
\texttt{continuous\_domain} for the real continuity domain of that expression.
The interval-building code in \texttt{solve\_univariate\_inequality} then tests
sample points between critical points and accepts whole intervals whose samples
satisfy the inequality.

The problematic rule reported in \texttt{continuous\_domain} treats rational
powers with odd denominator as real-valued over all real bases:

\begin{lstlisting}[style=bugpython]
if atom.exp.is_rational and den.is_odd:
    pass
\end{lstlisting}

That rule matches \texttt{real\_root}-style semantics, but it is too broad for
principal \texttt{Pow} semantics. For \(x^{1/3}-1\), it causes
\texttt{continuous\_domain} to return all real numbers rather than restricting
the expression to \(x \geq 0\). With the domain incorrectly set to
\(\mathbb{R}\), the inequality solver samples the interval \((-\infty,1)\);
the selected sample point satisfies the inequality, so the solver appends the
entire interval \((-\infty,1]\). No later filtering removes the negative inputs
where the principal fractional power is complex.

The suggested fix direction in the diagnosis is to make
\texttt{continuous\_domain} distinguish principal \texttt{Pow} from real-root
semantics. For rational non-integer powers of a real expression, a plausible
fix is to require the base to be nonnegative unless SymPy can prove the
principal value remains real on a larger set. The inequality solver should also
intersect interval results with the real-valued domain of the relational
expression.

\subsection{Additional Failing Instances}

The independent verification checks the representative case together with the
positive integer parameters \(m=1,\ldots,6\). In each case, SymPy returns an
interval of the form \((-\infty,m^3]\), which includes negative real inputs such
as \(-1\). The expected behavior is uniformly \([0,m^3]\), because the principal
power \(x^{1/3}\) is real-valued only for \(x\geq0\), and on that domain
\(x^{1/3}\leq m\) is equivalent to \(x\leq m^3\) for these positive \(m\).

\begin{center}
\small
\renewcommand{\arraystretch}{1.3}
\begin{tabular}{@{}>{\raggedright\arraybackslash}p{0.44\linewidth}>{\raggedright\arraybackslash}p{0.23\linewidth}>{\raggedright\arraybackslash}p{0.23\linewidth}@{}}
\toprule
\textbf{Input / Expression} & \textbf{SymPy behavior} & \textbf{Expected behavior} \\
\midrule
\(\texttt{solve\_univariate\_inequality}(x^{1/3}\leq 1)\) & \(\texttt{Interval(-oo, 1)}\) & \(\texttt{Interval(0, 1)}\) \\
\(\texttt{solve\_univariate\_inequality}(x^{1/3}\leq 2)\) & \(\texttt{Interval(-oo, 8)}\) & \(\texttt{Interval(0, 8)}\) \\
\(\texttt{solve\_univariate\_inequality}(x^{1/3}\leq 3)\) & \(\texttt{Interval(-oo, 27)}\) & \(\texttt{Interval(0, 27)}\) \\
\(\texttt{solve\_univariate\_inequality}(x^{1/3}\leq 4)\) & \(\texttt{Interval(-oo, 64)}\) & \(\texttt{Interval(0, 64)}\) \\
\(\texttt{solve\_univariate\_inequality}(x^{1/3}\leq 5)\) & \(\texttt{Interval(-oo, 125)}\) & \(\texttt{Interval(0, 125)}\) \\
\(\texttt{solve\_univariate\_inequality}(x^{1/3}\leq 6)\) & \(\texttt{Interval(-oo, 216)}\) & \(\texttt{Interval(0, 216)}\) \\
\bottomrule
\end{tabular}
\end{center}

\subsection{Related Issues Assessment}

The bug-finding harness performed a best-effort deduplication check and
classified this as a new instance of a known failure mode. The closest public
material identified was SymPy documentation distinguishing principal roots from
\texttt{real\_root}, together with broader public awareness that principal-root
domain handling can be subtle. The available evidence did not identify a clear
public duplicate for this exact \texttt{solve\_univariate\_inequality} input or
for the exact wrong interval \(\texttt{Interval(-oo, 1)}\). This case is
therefore individually actionable as a concrete regression test and issue,
even though the general semantic pitfall is known.

\subsection{Proposed SymPy Regression Test}

\medskip

\begin{lstlisting}[style=bugpython]
import pytest
from sympy import Interval, S, symbols
from sympy.solvers.inequalities import solve_univariate_inequality

@pytest.mark.parametrize("m", [1, 2, 3, 4, 5, 6])
def test_fractional_power_inequality_real_domain(m):
    x = symbols("x")
    assert solve_univariate_inequality(x**(S(1)/3) <= m, x, relational=False) == Interval(0, m**3)
\end{lstlisting}

\end{document}
